% =============================================================
% chapters/vysnovky.tex — ЗАГАЛЬНІ ВИСНОВКИ
% Висновки до~всієї роботи (підсумок Розділу 1 + Розділу 2).
% Висновки до~окремих розділів подаються наприкінці кожного rozdil.
% TODO[T1.7, T2.10]: розширити після завершення обох розділів.
% =============================================================

\section*{\centering ВИСНОВКИ}\addcontentsline{toc}{section}{ВИСНОВКИ}

У~кваліфікаційній роботі виконано бібліометричне обґрунтування та~спроєктовано інформаційно-аналітичну систему моніторингу освітньої діяльності закладів професійної освіти на~основі 1998~публікацій бази Web of~Science Core Collection (1997--2025) і~нормативної бази Закону України «Про професійну освіту»~№~4574-IX.

\textbf{За~результатами дослідження отримано такі основні висновки:}

\begin{enumerate}
    \item Описова статистика (підрозділ~\ref{sec:p1-descriptive}) показала стале зростання кількості публікацій із~моніторингу професійної освіти з~піком у~2019--2023~роках. Географічно дослідження зосереджені у~Європі (Німеччина, Нідерланди, Велика Британія), що~відображає розвиненість систем дуальної освіти в~цих країнах. Провідними науковими джерелами є~спеціалізовані журнали з~професійної освіти.

    \item Бібліометричний аналіз мережі ключових слів у~VOSViewer (підрозділ~\ref{sec:p1-clusters}) виявив 4~тематичні кластери: «Системи та~політика ПТО» (11~термінів), «Результати навчання» (8), «Педагогічна практика» (6), «Когнітивні результати» (3). Центральними термінами є~\textit{vet} (TLS\,=\,17\,857), \textit{study} (14\,924), \textit{student} (14\,844), \textit{training} (13\,158). Динаміка поля свідчить про~зростання уваги до~ролі роботодавців (\textit{company}, 2019,85) та~контекстуалізованого навчання (\textit{context}, 2019,09).

    \item Порівняльний аналіз (підрозділ~\ref{sec:p1-comparison}) продемонстрував високу відтворюваність кількісних метрик: кореляція частот входжень $r=0{,}923$, Total~Link~Strength $r=0{,}919$, середнього року публікації $r=0{,}74$, середньої цитованості $r=0{,}74$, нормалізованої цитованості $r=0{,}79$.

    \item Кластеризація алгоритмом Лувена в~Python оптимально дає 3~кластери (ARI\,=\,0,47; NMI\,=\,0,61), об'єднуючи кластери «Педагогічна практика» та~«Когнітивні результати» VOSViewer. Це~пояснюється майже повною зв'язністю мережі (378/378~пар) та~відмінностями в~алгоритмах кластеризації.

    \item Картування 4~бібліометричних кластерів у~7~блоків Specifikatsiya~v2.0 (підрозділ~\ref{sec:p1-impl-mapping}) і~5~дизайн-вимог для~переміщених закладів П(ПТ)О Луганщини (підрозділ~\ref{sec:p1-impl-displaced}) забезпечують безпосередню прив'язку глобальних висновків до~реформи професійної освіти України, рухомої трьома драйверами: воєнним часом, повоєнним відновленням і~гармонізацією з~ЄС у~рамках EU4Skills, EQAVET та~ETF reporting.

    \item Розроблений у~Розділі~2 артефакт реалізує всі~24~форми Specifikatsiya~v2.0 у~8~функціональних модулях з~ідемпотентним прийомом даних, декларативною ABAC-моделлю доступу за~JSON-AST граматикою з~оператором \texttt{territorial\_contains}, SHA-256 хеш-ланцюгом аудиту та~підтримкою переміщених закладів через~розрізнення юридичної/фактичної адреси. Емпірична евалюація через~синтетичні бенчмарки і~експертне оцінювання $n\geq 3$ підтверджує функціональну повноту, ABAC-коректність ($\sim$160~тестів no-bypass), цілісність аудит-ланцюга на~10\,000~upsert та~відповідність WCAG~2.1~AA.
\end{enumerate}

\textbf{Обмеження дослідження:}
\begin{itemize}
    \item використання лише однієї бібліографічної бази (Web of~Science);
    \item обмеження англомовними публікаціями;
    \item різний обсяг вибірок для~VOSViewer (1000~записів) та~Python (1998~записів), що~впливає на~абсолютні значення метрик;
    \item текстовий аналіз (text~mining) охоплює лише заголовки та~абстракти, а~не повні тексти;
    \item емпірична евалюація системи (Розділ~2) обмежена синтетичними бенчмарками та~експертним оцінюванням ($n\geq 3$) без~повномасштабного пілотного впровадження.
\end{itemize}

\textbf{Перспективи подальших досліджень:}
\begin{itemize}
    \item розширення бібліометричного аналізу на~бази Scopus та~ERIC для~отримання повнішої картини;
    \item включення україномовних та~інших неангломовних публікацій;
    \item повномасштабне пілотне впровадження системи в~одному обласному центрі моніторингу ПТО з~подальшим лонгітюдним оцінюванням;
    \item інтеграція з~ЄДЕБО, додавання КЕП-підпису офіційних звітів, запровадження ML-передбачення відсіву здобувачів;
    \item адаптація під~ETF reporting та~гармонізацію з~EQAVET у~рамках EU4Skills.
\end{itemize}
